\documentclass{article}
\input{../lib.sty}

\title{\texttt{fn discrete\_laplacian\_scale\_to\_accuracy}}
\author{\MichaelShoemateAuthor}

\begin{document}
\maketitle
\contrib

This document contains materials associated with \rustdoc{accuracy/fn}{discrete\_laplacian\_scale\_to\_accuracy}.

\begin{definition}
    Let $z$ be the true value of the statistic and $X$ be the random variable the noisy release is drawn from.
    Define $Y = |X - z|$, the distribution of DP errors.
    For a statistical significance level $\alpha\in(0,1]$, a positive accuracy threshold $a$ is sufficient when
    \begin{equation}
        P[Y \ge a]\leq\alpha
    \end{equation}
\end{definition}

\begin{theorem}
    For positive finite scale $s$ and $X\sim\mathcal{L}_\mathbb{Z}(z,s)$, a sufficient real-valued accuracy threshold is
    \begin{equation}
        a = s \cdot ln(2 / (\alpha (e^{1/s} + 1))) + 1
    \end{equation}
\end{theorem}

\begin{proof}

Consider that the distribution of $(X - z) \sim \mathcal{L}_\mathbb{Z}(0, s)$.
Then the PMF of $Y$ is:
\begin{equation}
    \forall y \in \mathbb{Z}_{\ge 0} \qquad g(y) = (1 + 1[y \ne 0]) \frac{1 - e^{-1/s}}{1 + e^{-1/s}} e^{-y/s}
\end{equation}

The purpose of the indicator function is to avoid double-counting zero.

For positive integers $a$, derive the tail probability:
\begin{align*}
    P[Y \ge a] &= 1 - P[Y < a] \\
    &= 1 - \sum_{y=0}^{a - 1} g(y) && \text{where $g(y)$ is the distribution of Y} \\
    &= 1 - \frac{1 - e^{-1/s}}{1 + e^{-1/s}} \left(1 + 2\sum_{y=1}^{a - 1} e^{-y/s}\right) \\
    &= 2 \frac{e^{(1 - a)/s}}{e^{1/s} + 1}
\end{align*}

Extend the tail expression to positive real arguments and solve for the crossing at level $\alpha$:
\begin{align*}
    2 \frac{e^{(1 - a)/s}}{e^{1/s} + 1} &= \alpha \\
    e^{(1 - a)/s} &= \alpha (e^{1/s} + 1) / 2 \\
    a &= 1 - s \cdot ln(\alpha (e^{1/s} + 1) / 2) \\
    a &= s \cdot ln(2 / (\alpha (e^{1/s} + 1))) + 1
\end{align*}

The expression is decreasing in $a$. Since $Y$ is integer-valued, the actual tail at a positive real threshold equals the tail at its ceiling, so the threshold is sufficient in exact arithmetic.
At scale zero the noise is identically zero and the implementation returns zero accuracy.
The floating-point implementation evaluates the expression numerically; this derivation does not establish a directed-rounding guarantee.
\end{proof}

\end{document}
